%!TEX root = main.tex
%!TEX spellcheck = en_US

\section{Introduction}
\label{sec:intro}

Ensuring the stability, security, and performance of network infrastructures hinges on careful management when introducing a new routing device or modifying the configuration of an existing device. Even a single device's misconfiguration can lead to black holes, unintended network access, inefficient routes, and a range of other performance and security vulnerabilities. For instance, a misconfigured access control list (ACL) may unintentionally allow unauthorized traffic, exposing the device (and potentially the broader network) to malicious attacks; a buggy routing policy may create loops, thereby degrading performance and rendering local services inaccessible.

The most accurate but costly approach to verifying network configuration updates is to involve a domain expert, such as a network engineer. These experts bring an extensive and comprehensive set of background knowledge gathered over years of experience, vast documentation, and an understanding of standard or best practices in routing configurations. Given a configuration under review, the expert iterates through it, identifies issues, and formulates these issues into policies that future model checkers can use to detect similar misconfigurations. However, this manual process is labor-intensive, costly, and hard to scale, especially in large and dynamic network environments.

To address these limitations, researchers and practitioners have developed automated tools such as model checkers and consistency checkers.
- \textbf{Model checkers}~\cite{fogel2015general, beckett2017general, abhashkumar2020tiramisu, prabhu2020plankton, zhang2022sre, steffen2020netdice, ye2020hoyan, ritchey2000using,al2011configchecker, jeffrey2009model} rely on manually defined \emph{policies}, which include correctness constraints and protocol-specific requirements that configurations must satisfy, enabling them to detect complex misconfigurations accurately. However, curating these policies is challenging. Tailoring rules to different scenarios, device types, software versions, and vendor-specific features requires significant effort, and creating a comprehensive set of policies is often infeasible. Unlike human experts, model checkers cannot relate disparate pieces of information within a configuration without explicit rules.
- \textbf{Consistency checkers}~\cite{kakarla2024diffy,kakarla2020finding, le2006minerals, feamster2005detecting, tang2021campion, le2008detecting, le2006characterization}, on the other hand, automate the process by comparing configurations across devices or best-practice baselines. These tools flag deviations as potential issues. While this approach requires minimal manual intervention, it lacks semantic depth. Deviations do not necessarily imply faults, and consistency-based tools cannot leverage domain knowledge or policies derived from expert insights.

Given these challenges, recent efforts start to explore the feasibility of using large language models (LLMs) as domain experts to carry out the task of misconfiguration detection. LLMs are pre-trained on extensive corpora that encompass network-related texts, including protocol specifications, vendor documentation, and publicly available configurations. This makes them analogous to human experts with broad domain knowledge. By feeding the configuration under review to an LLM, the model can iterate through it and identify issues, potentially offering a scalable and automated alternative to manual inspection.

Existing LLM-based tools~\cite{bogdanov2024leveraging,chen2024automatic,wang2024identifying,liu2024large, wang2024netconfeval, lian2023configuration} attempt to achieve this by passing entire or partitioned configurations to the LLM in prompts. However, this approach has limitations. Like human engineers, LLMs require \emph{context} to identify issues. Human experts do not examine configurations in a rigid, top-down fashion; instead, they identify relevant configuration lines and dynamically ``jump'' to related parts of the configuration, regardless of partition boundaries. Similarly, LLMs perform better when provided with concise, relevant context that reflects these interactions. Feeding the entire configuration at once risks exceeding token limits, and splitting configurations into arbitrary partitions can obscure dependencies between different parts. Furthermore, overloading the model with irrelevant or excessive information can dilute its reasoning capabilities, much like a human struggling to recall long-past details.

In this paper, we propose a framework for efficiently extracting and feeding context into LLMs to address these issues and improve misconfiguration detection in routing devices. Specifically, we introduce the \emph{Context-Aware Iterative Prompting} (\sysname{}) framework, which tackles three key challenges:

\mypara{Challenge 1: Automatic and Accurate Context Mining} Router configurations are long and complex~\cite{benson2009complexitymetrics}, featuring interdependent lines.
Imagine a domain expert tasked with analyzing a configuration file. Instead of reviewing every line exhaustively, they focus on key details and related information. 
Similarly, automatically extracting and selecting relevant context is essential to reduce computational expense and improve detection accuracy in LLMs. \textbf{Solution:} We model the configuration as a hierarchical tree, where each line corresponds to a path from the root to a parameter value. This structure enables systematic identification of three main categories of context---\emph{neighboring}, \emph{similar}, and \emph{referenced} configuration lines—ensuring the LLM receives a concise yet sufficient corpus for analysis.

\mypara{Challenge 2: Parameter-Value Ambiguity in Referenceable Context} Configurations contain pre-defined and user-defined parameters, with the latter often requiring specialized context for accurate interpretation.
Consider a scenario where an expert encounters a user-defined parameter (\eg, an ACL name) and must track its references across the configuration. Similarly, LLMs require clarity in resolving such references.
Mislabeling or omitting these references inflates computational costs and increases error risk. \textbf{Solution:} We employ an existence-based labeling approach within the configuration tree to distinguish user-defined values by detecting their occurrences as internal nodes in other paths. These classifications are refined with a majority-voting scheme to ensure consistency across contexts.

\mypara{Challenge 3: Managing Context Overload in Prompts} 
When faced with an overwhelming volume of irrelevant information, even the best experts can lose focus. Likewise, even after extracting relevant information, the quantity of required context can exceed a single prompt's capacity. Overloading an LLM with excessive detail reduces coherence and degrades performance. \textbf{Solution:} In our iterative prompting framework, the LLM explicitly requests further context when needed, incrementally refining the prompt and focusing on the most critical lines. This approach mitigates overload and enables more precise detection.

We evaluate \sysname{} through two distinct case studies. First, we analyze real-world device configuration snapshots, inject realistic updates (including both correct modifications and latent misconfigurations), and demonstrate that \sysname{} outperforms partition-based LLM approaches, model checkers, and consistency checkers by at least 30\% in error detection accuracy. Second, we exhaustively evaluate \sysname{} on configuration snapshots from a medium-sized campus network. \sysname{} identifies over 20 previously overlooked misconfigurations, highlighting its practical utility.

